\section{An extension of two stable bundles}\label{sec:two-factor}

Let $(X,\omega)$ again be a compact connected K\"ahler manifold. Let
$F,G$ be nonisomorphic stable bundles of equal slope, with ranks $a,b$,
and put $r=a+b$. Fix Hermitian--Einstein metrics $h_F,h_G$, whose common
Einstein constant is $\kappa=2\pi\mu_\omega(F)/V$, and a nonzero harmonic
form
\[
 \beta\in\Omega^{0,1}(\Hom(G,F)),\qquad
 \db\beta=0,\qquad \db^*\beta=0.
\]
All inner products in this section use $h_0=h_F\oplus h_G$ and the fixed
K\"ahler form. We use the convention
$\langle u,v\rangle=\int_X\tr(u^*v)\dvol$.

On the fixed smooth bundle $F\oplus G$ define
\begin{equation}\label{eq:two-family}
 N=\begin{pmatrix}0&\beta\\0&0\end{pmatrix},\qquad
 \db_t=\db_0+tN,\qquad t\in\mathbb R.
\end{equation}
This is integrable in every complex dimension because $\db_0N=0$ and
$N\wedge N=0$. Denote the resulting holomorphic bundle by $E_t$.
On trace-free endomorphisms put
\[
 D_t=D_0+tB,\qquad D_0=\db_{0,\End},\qquad
 Bu=[N,u],\qquad \Delta_t=D_t^*D_t.
\]
These nonnegative self-adjoint elliptic operators have the fixed domain
$H^2$ in $L^2$. Eigenvalues count complex multiplicity and include zero.
Let $P_F$ be the orthogonal projection onto $F$ and set
\begin{equation}\label{eq:two-q}
 q=P_F-\frac ar\id
 =\diag\left(\frac br\id_F,-\frac ar\id_G\right),\qquad
 V_q=\lVert q\rVert_{L^2}^2=\frac{Vab}{r}.
\end{equation}

Write $\beta=\sum_\nu\beta_\nu\bar\theta^\nu$ in a local unitary
coframe, and define the global self-adjoint trace-free endomorphism
\begin{equation}\label{eq:two-C}
 C=\diag\left(\sum_\nu\beta_\nu\beta_\nu^*,
                  -\sum_\nu\beta_\nu^*\beta_\nu\right).
\end{equation}
Stability gives $\Ker\Delta_0=\CC q$, as verified below. Let $P$ be the
$L^2$ projection onto this kernel, $Q=\id-P$, and
$R=(\Delta_0|_{QH^2})^{-1}:QL^2\to QH^2$. Define
\begin{equation}\label{eq:two-coefficients}
 c=\frac{r\lVert\beta\rVert_{L^2}^2}{Vab},\qquad
 z=C-cq=QC,\qquad d=\frac{\langle z,Rz\rangle}{V_q}.
\end{equation}

\begin{theorem}\label{thm:two-factor}
For the family \eqref{eq:two-family}, $E_t$ is simple and strictly
semistable for every $t\ne0$. Exactly one eigenvalue of $\Delta_t$ tends
to zero, and
\begin{equation}\label{eq:two-expansion}
 \lambda_1(t)=ct^2-dt^4+O(t^6),\qquad c>0,\quad d\ge0.
\end{equation}
There are $\delta,t_0>0$ such that $\lambda_2(t)\ge\delta$ for
$|t|<t_0$. The normalized first eigensection, with phase chosen so that
$\langle q,u_t\rangle>0$, satisfies for every $k\ge0$
\begin{equation}\label{eq:two-section}
 u_t=\frac{q}{\sqrt{V_q}}-\frac{t^2}{\sqrt{V_q}}Rz+O_{C^k}(t^3).
\end{equation}
The mean curvature and its averaged trace-free norm are exactly
\begin{equation}\label{eq:two-curvature}
 K_t=\kappa\id+t^2C,\qquad
 \varepsilon_1(t)=\frac{t^2}{V}\int_X\lVert C\rVert_{\op}\dvol.
\end{equation}
Moreover, $\lambda_1(t)\le ct^2$ for every real $t$.
\end{theorem}

\begin{proof}
Nonisomorphism and stability at equal slope imply
\[
 H^0(\Hom(F,G))=H^0(\Hom(G,F))=0,\qquad
 H^0(\End F)=\CC\id_F,\quad H^0(\End G)=\CC\id_G.
\]
Thus the trace-free holomorphic endomorphisms of the split bundle are
precisely $\CC q$. This proves the assertion about $\Ker\Delta_0$ and
the existence of $R$.

For $t\ne0$, an endomorphism of $E_t$ preserves $F$, since its composite
$F\to E_t\to G$ vanishes. It induces scalars $\alpha,\eta$ on $F,G$,
and its upper-right block $v$ satisfies
\[
 \db v+t(\eta-\alpha)\beta=0.
\]
Pairing with $\beta$ gives $\eta=\alpha$. The block $v$ is then
holomorphic and hence zero, proving simplicity. For any subsheaf of
$E_t$, its intersection with $F$ and image in $G$ have slopes at most
the common slope whenever their ranks are nonzero. Degree and rank
additivity prove semistability, and the equal-slope subbundle $F$
rules out stability.

Expand on the common domain:
\begin{equation}\label{eq:two-operator}
 \Delta_t=\Delta_0+tA_1+t^2A_2,\qquad
 A_1=D_0^*B+B^*D_0,\quad A_2=B^*B.
\end{equation}
Since $D_0q=0$, $Bq=-N$, and $D_0^*N=0$ by harmonicity, one has
$A_1q=0$. The adjoint formula
$B^*v=\sum_\nu[N_\nu^*,v_\nu]$ gives
\begin{equation}\label{eq:two-action}
 A_2q=C,\qquad \Delta_tq=t^2C,\qquad
 \langle q,C\rangle=\lVert\beta\rVert_{L^2}^2.
\end{equation}
In particular $PC=cq$, and the Rayleigh quotient of $q$ gives
$\lambda_1(t)\le ct^2$.

The zero eigenvalue of $\Delta_0$ is isolated and simple. The analytic
family \eqref{eq:two-operator} has compact resolvent and constant domain,
so its spectral projection near zero is analytic and of rank one for
small $|t|$ \cite[Chapter~VII]{Kato1995}. The remaining spectrum stays
bounded away from zero. For clarity, on a contour separating zero from
the positive spectrum, $(\Delta_0-\zeta)^{-1}$ maps $L^2$ to $H^2$,
and
$(tA_1+t^2A_2)(\Delta_0-\zeta)^{-1}$ has norm $O(|t|)$ on $L^2$.
A Neumann series constructs the perturbed resolvent and its contour
integral. The same argument between $H^s$ and $H^{s+2}$ gives Taylor
expansions of the eigensection in every $C^k$ norm.

Choose the analytic eigensection $v(t)$ with
$\langle q,v(t)\rangle=V_q$ and write
\[
 v(t)=q+tv_1+t^2v_2+t^3v_3+\cdots,\qquad v_j\perp q.
\]
The order-$t$ equation gives $\lambda'(0)=0$ and $v_1=0$. At order
$t^2$, projection onto $q$ and $q^\perp$ gives respectively the
coefficient $c$ and
\[
 \Delta_0v_2+C=cq,\qquad v_2=-Rz.
\]
The order-$t^3$ coefficient vanishes because $A_1q=0$. If
$\lambda_{[4]}$ denotes the coefficient of $t^4$, then
\[
 V_q\lambda_{[4]}
 =\langle q,A_1v_3+A_2v_2\rangle
 =-\langle z,Rz\rangle.
\]
This coefficient is $-d$, with $d\ge0$ because $R$ is positive.
Conjugation by $J=\diag(-\id_F,\id_G)$ is unitary and intertwines
$D_t$ with $D_{-t}$. Uniqueness of the isolated branch makes its
eigenvalue even in $t$, proving the $O(t^6)$ remainder in
\eqref{eq:two-expansion}. Finally,
$\lVert v(t)\rVert^2=V_q+O(t^4)$, so normalization gives
\eqref{eq:two-section}.

The Chern connection is $\nabla_t=\nabla_0+t(N-N^*)$, where the adjoint
also changes the form type. The K\"ahler identity
$\db^*\beta=-\sqrt{-1}\Lambda_\omega\partial\beta=0$ and its
adjoint kill the linear term in the contracted curvature. In a unitary
coframe with $\omega=\sqrt{-1}\sum\theta^\nu\wedge\bar\theta^\nu$,
the quadratic term is
\[
 \sqrt{-1}\Lambda_\omega(N-N^*)\wedge(N-N^*)=C.
\]
Its trace is zero, proving \eqref{eq:two-curvature}.
\end{proof}

To view this degeneration on the fixed holomorphic bundle $E_1$, set
$g_t=\diag(t\id_F,\id_G)$ for $t>0$. Then
\begin{equation}\label{eq:two-gauge}
 \db_t=g_t\db_1g_t^{-1},\qquad
 h_t=g_t^*h_0=t^2h_F\oplus h_G.
\end{equation}
The map $g_t:(E_1,h_t)\to(E_t,h_0)$ is a holomorphic isometry and
identifies the endomorphism Laplacians unitarily. Thus $h_t$ is an
approximate Hermitian--Einstein family on $E_1$. The convergence in
\eqref{eq:two-section} is in the gauge $g_t$ of \eqref{eq:two-gauge}. Its limit $u_0$
has two distinct fiberwise eigenvalues, and
\[
 P_F=\sqrt{V_q}\,u_0+\frac ar\id,\qquad P_G=\id-P_F.
\]

